% =============================================================
%  后记
% =============================================================
\ctexset{section/numbering=false}

\Chapter{后记}

\vspace{0.4cm}

写完这份执行手册，感觉有点人格分裂。

写的过程先要摆脱“豆角爸爸”这个当事人的身份，把自己变成小豆角学校的老师，同时又是小豆角本人——
先以“老师”的身份教会“学生”，以“学生”的身份判断老师讲的是否足够简单易懂，以此判断自己讲的东西
是否被自己理解，形成逻辑闭环，然后从“学生”切换成“豆角爸爸”，再以“豆角爸爸”的身份来模拟跟小豆
角互动，以此来指导遣词造句细节。

一会儿站在岸上，一会儿泡在水里，逻辑难免混乱。

不过整书复核、校对下来，勉强能用了。先用着吧。

\vfill

\begin{center}
{\kaishu\large 豆角爸爸}\\[4pt]
{\small\color{tipc} https://www.bd7iwd.com}
\end{center}

\vspace{0.6cm}
